\section{Méthodologie}

Cette section décrit la méthodologie expérimentale employée pour évaluer
les performances du système, incluant le protocole de benchmark des modèles
d'embedding, l'évaluation des stratégies de recherche, et la métrique de
qualité des réponses générées.

\subsection{Benchmark des Modèles d'Embedding}

\subsubsection{Corpus de Test}

Un corpus de 26 articles juridiques marocains a été constitué, couvrant cinq
domaines : pénal (7 articles, n° 147--436), trafic (5 articles, n° 172--180),
familial (5 articles, n° 48--78), civil (4 articles, n° 1--100), et
commercial (4 articles, n° 1--55). Chaque article a été soigneusement
sélectionné pour inclure des articles proches (hard negatives) — par exemple,
l'article 147 (vol simple) et l'article 148 (vol aggravé) du Code Pénal —
afin de tester la capacité des modèles à discriminer des textes
juridiquement distincts mais sémantiquement similaires.

\subsubsection{Variantes de Requêtes}

Pour chaque groupe d'articles, cinq à sept variantes de requêtes ont été
générées :

\begin{itemize}
    \item \textbf{Arabe formel} : question en arabe standard moderne
    (\textarabic{\texttt{ما هي عقوبة السرقة البسيطة؟}})
    \item \textbf{Darija marocain} : même question en dialecte marocain
    (\textarabic{\texttt{شنو هي العقوبة ديال السرقة؟}})
    \item \textbf{Court (keywords)} : 2-3 mots-clés
    (\textarabic{\texttt{عقوبة السرقة}})
    \item \textbf{Long (descriptif)} : phrase descriptive détaillée
    (\textarabic{\texttt{أريد معرفة العقوبة المقررة قانونا لجريمة السرقة البسيطة وفق
    القانون الجنائي المغربي}})
    \item \textbf{Cross-lingue français} : question en français
    (\texttt{Quelle est la peine pour le vol simple au Maroc?})
    \item \textbf{Cross-lingue anglais} : question en anglais
    (\texttt{What is the penalty for simple theft in Morocco?})
\end{itemize}

Au total, 65 variantes de requêtes ont été évaluées.

\subsubsection{Modèles Évalués}

Cinq modèles d'embedding multilingues ont été comparés :

\begin{table}[H]
\centering
\caption{Modèles d'embedding évalués}
\resizebox{\columnwidth}{!}{%
\begin{tabular}{llc}
\toprule
\textbf{Modèle} & \textbf{Identifiant} & \textbf{Dim.} \\
\midrule
E5-large & \texttt{intfloat/multilingual-e5-large} & 1024 \\
BGE-M3   & \texttt{BAAI/bge-m3}                 & 1024 \\
LaBSE    & \texttt{sentence-transformers/LaBSE}  & 768  \\
MPNet    & \texttt{paraphrase-multilingual-mpnet-base-v2} & 768 \\
E5-base  & \texttt{intfloat/multilingual-e5-base} & 768  \\
\bottomrule
\end{tabular}
}
\end{table}

Tous les modèles ont été chargés via \texttt{sentence-transformers} avec
normalisation L2 des embeddings. L'évaluation a été effectuée sur GPU
lorsque disponible.

\subsubsection{Métriques}

Pour chaque paire (modèle, variante), la similarité cosinus entre
l'embedding de la requête et les embeddings des 26 articles a été calculée.
Les métriques suivantes ont été agrégées sur tous les groupes de requêtes :

\begin{itemize}
    \item \textbf{Précision@k} ($k \in \{1,3,5\}$) : proportion de requêtes
    où l'article pertinent est dans les $k$ premiers résultats
    \item \textbf{MRR} (Mean Reciprocal Rank) : moyenne du rang inverse de
    l'article pertinent
    \item \textbf{NDCG@5} (Normalized Discounted Cumulative Gain) : qualité
    du classement pondérée par la position
    \item \textbf{Précision de domaine} : proportion de requêtes où le
    meilleur résultat appartient au domaine juridique correct
    \item \textbf{Hit@10} : proportion de requêtes où l'article pertinent
    est dans les 10 premiers résultats
\end{itemize}

\subsection{Benchmark des Stratégies de Recherche}

\subsubsection{Configuration Expérimentale}

Un benchmark a été conçu pour comparer les trois stratégies de recherche
implémentées : dense (pgvector cosine), BM25 (rank-bm25 avec tokenizer
arabe), et RRF (fusion des deux). Le jeu de test comprend 32 questions
juridiques couvrant 7 domaines (trafic, pénal, familial, civil, commercial,
violations, entreprises), avec un top\_k=5.

\subsubsection{Métriques de Recherche}

\begin{itemize}
    \item \textbf{Precision@5} : proportion d'articles pertinents parmi
    les 5 résultats
    \item \textbf{Recall@5} : proportion d'articles pertinents retrouvés
    \item \textbf{MRR@5} : Mean Reciprocal Rank
    \item \textbf{NDCG@5} : Discounted Cumulative Gain normalisé
    \item \textbf{Hit@5} : proportion de questions avec au moins un
    article pertinent dans les résultats
\end{itemize}

\subsubsection{Métriques de Génération}

Pour chaque stratégie, les articles récupérés ont été passés au LLM (Gemma 4
2B, température 0,3, max\_tokens=512) pour générer une réponse. Cinq métriques
ont été évaluées :

\begin{itemize}
    \item \textbf{Answer Relevance} : similarité cosinus entre les embeddings
    de la question et de la réponse (0--1)
    \item \textbf{Faithfulness} : proportion de phrases de la réponse
    ancrables dans les documents retrouvés par chevauchement lexical (0--1)
    \item \textbf{Answer Correctness} : score composite ROUGE-L + BERTScore
    F1 (0--1)
\end{itemize}

\subsubsection{Analyse par Domaine et Niveau de Difficulté}

Les 32 questions du benchmark RAG complet ont été réparties sur 7 domaines
juridiques (pénal, trafic, familial, civil, commercial, fiscal, judiciaire)
et classées en trois niveaux de difficulté : \textit{easy} (10 questions,
vocabulaire général, un article attendu), \textit{medium} (11 questions,
deux articles attendus), et \textit{hard} (11 questions, vocabulaire technique
spécifique, articles adjacents à discriminer). Cette stratification permet
d'identifier les faiblesses spécifiques du système par domaine et par type
de requête.

\subsubsection{Significativité Statistique}

Pour le benchmark d'embedding, des tests $t$ de Student appariés ont été
appliqués entre chaque paire de modèles afin de déterminer si les différences
observées sont statistiquement significatives ($\alpha = 0,05$).

\subsection{Implémentation et Reproductibilité}

L'ensemble du système est conteneurisé via Docker Compose. Les services
suivants sont définis :

\begin{itemize}
    \item \textbf{PostgreSQL + pgvector} : base de données vectorielle
    avec pgAdmin pour la gestion
    \item \textbf{Service LLM} : Gemma 4 2B via llama.cpp, exposé sur
    le port 8001
    \item \textbf{API FastAPI} : application principale (port 8000)
    \item \textbf{Frontend SvelteKit} : interface utilisateur (port 5173)
\end{itemize}

Les benchmarks d'embedding et de recherche sont reproduisibles via les
notebooks Jupyter dans le dossier \texttt{notebooks/}, et les résultats
sont stockés dans \texttt{benchmark/output/}.
